%%%%%%%%%%%%%%%%%%%%%%%%%%%%%%%%%
% Introduction
%%%%%%%%%%%%%%%%%%%%%%%%%%%%%%%%%
\section{Introduction}
\label{sec:intro_ent}

The Bernoulli set model~\cite{bernoulliSets} provides a compositional framework
for random approximate sets---data structures whose membership queries may err
with known false positive and false negative rates.
A natural information-theoretic question is: how much uncertainty do these
errors carry?

\paragraph{Motivation.}
The joint entropy of error counts quantifies the inherent unpredictability of a
Bernoulli set realization.
This entropy has direct implications for coding: any lossless representation of
the error pattern requires at least this many bits on average.
The space complexity question---how many bits are needed to implement an
approximate set---connects the abstract Bernoulli model to the concrete cost of
probabilistic data structures.

\paragraph{Contribution.}
We make two contributions.
First, we derive the joint entropy of the false positive and false negative
counts (\cref{thm:joint_entropy_fp_fn}), showing that it decomposes additively
by element independence and admits the closed-form asymptotic expansion
\begin{equation}
    \log_2(2\pi e)
    + \frac{1}{2}\log_2\!\bigl(
        (u-m) \cdot m \cdot \fprate(1-\fprate) \cdot \fnrate(1-\fnrate)
      \bigr)
    + \mathcal{O}\!\left(\frac{u}{m(u-m)}\right).
\end{equation}
We extend this to the case where the number of positives is itself uncertain
(\cref{sec:entropy:pos_neg}).
Second, we restate and prove the classical space lower bound of Carter et
al.~\cite{carterExact} for positive approximate sets
(\cref{pst:approx_l_b}): any implementation over a countably infinite universe
requires at least $-\log_2 \fprate$ bits per element.
We compare this to the Bloom filter's $-(1/\ln 2)\log_2\fprate$ bits per
element and note that the general case $0 < \tprate < 1$ remains open
(\cref{rem:general_lower_bound}).

\paragraph{Related work.}
Shannon~\cite{shannonBSC} introduced the entropy of binary channels; Cover and
Thomas~\cite{coverThomas} provide the binomial entropy expansion we use.
Carter et al.~\cite{carterExact} established the space lower bound.
Bloom~\cite{bf} introduced the eponymous filter.
The Bernoulli set model and its compositional algebra are developed in a
companion paper~\cite{bernoulliSets}.

\paragraph{Organization.}
\Cref{sec:prelim_ent} restates the Bernoulli axioms.
\Cref{sec:entropy} derives the joint entropy.
\Cref{sec:space_comp} establishes the space complexity lower bound.
\Cref{sec:conclusion_ent} concludes with open problems.
